% GENERATED FILE. Edit canonical lessons, then run npm run generate:books.
% canonical-source-hash: fnv1a64:036deb7cc9c185cf
% canonical-lessons: UR-C530-marzi, UR-C530-razi, UR-C530-ikhtilaf, UR-C530-mustarad, UR-C530-palatna

\chapter{Yes or No}
\label{ch:ur-yes-or-no}
\hlchaptermodality{530}

\begin{chapteropening}
\textbf{By the end of this chapter:} \emph{I can talk about agreeing and turning things down.}

Complete the last lesson of chapter 530: I can talk about agreeing and turning things down.
\end{chapteropening}

\section[\texorpdfstring{\ur{مرضی} (marzī)}{مرضی (marzī)}]{\ur{مرضی} (marzī) — one's own choice; consent}

\label{lesson:UR-C530-marzi}



\begin{quote}
Before the new one: say the Urdu for complicated, then the Urdu for to tangle.
\end{quote}

\subsection*{You'll want to know: \ur{مرضی}}
\textbf{\ur{مرضی}} — \emph{marzī} — \textquotedblleft{}one's own choice; consent\textquotedblright{}.

\begin{quote}
\ur{اگر} \ur{آپ} \ur{کی} \ur{مرضی} \ur{ہے} \ur{تو} \ur{ہم} \ur{چلیں} \ur{گے۔}
\end{quote}

\begin{quote}
\emph{agar āp kī marzī hai to ham chaleṅ ge.}
\end{quote}

\begin{quote}
If you wish, we will go.
\end{quote}

Feminine.

\begin{cousinweb}
\textbf{\ur{مرضی}} (\emph{marzī}) is usually traced, through Persian, to Arabic \emph{marḍī}, \textquotedblleft{}pleasing, approved\textquotedblright{}.
\end{cousinweb}

\begin{grammarlens}[title={What you've built}]
One more word for saying what happens if.
\end{grammarlens}

\subsection*{Your turn}
\begin{itemize}
\raggedright
  \item \emph{Say it:} \emph{marzī}
  \item \emph{Say it:} \emph{marzī}, once more
  \item \emph{Say it:} \emph{agar āp kī marzī hai to ham chaleṅ ge.} — If you wish, we will go.
  \item \emph{Recall:} say the Urdu for complicated, then the Urdu for to tangle, then say \emph{marzī} again
\end{itemize}

\begin{tcolorbox}[breakable,colback=teal!4,colframe=teal!35!black,title={Before you move on}]
What does \ur{مرضی} mean? (\textquotedblleft{}One's own choice; consent\textquotedblright{}.) Say it once more.
\end{tcolorbox}
\section[\texorpdfstring{\ur{راضی} (rāzī)}{راضی (rāzī)}]{\ur{راضی} (rāzī) — willing; agreed}

\label{lesson:UR-C530-razi}



\begin{quote}
Before the new one: say the Urdu for to tangle, then the Urdu for one's own choice.
\end{quote}

\subsection*{You'll want to know: \ur{راضی}}
\textbf{\ur{راضی}} — \emph{rāzī} — \textquotedblleft{}willing; agreed\textquotedblright{}.

\begin{quote}
\ur{اگر} \ur{ابو} \ur{راضی} \ur{ہوئے} \ur{تو} \ur{میں} \ur{جاؤں} \ur{گا۔}
\end{quote}

\begin{quote}
\emph{agar abbū rāzī hue to maiṅ jāūṅ gā.}
\end{quote}

\begin{quote}
If my father agrees, I will go.
\end{quote}

\begin{cousinweb}
\textbf{\ur{راضی}} (\emph{rāzī}) is usually traced, through Persian, to Arabic \emph{rāḍī}, \textquotedblleft{}content, pleased\textquotedblright{}.
\end{cousinweb}

\begin{grammarlens}[title={What you've built}]
One more word for saying what happens if.
\end{grammarlens}

\subsection*{Your turn}
\begin{itemize}
\raggedright
  \item \emph{Say it:} \emph{rāzī}
  \item \emph{Say it:} \emph{rāzī}, once more
  \item \emph{Say it:} \emph{agar abbū rāzī hue to maiṅ jāūṅ gā.} — If my father agrees, I will go.
  \item \emph{Recall:} say the Urdu for to tangle, then the Urdu for one's own choice, then say \emph{rāzī} again
\end{itemize}

\begin{tcolorbox}[breakable,colback=teal!4,colframe=teal!35!black,title={Before you move on}]
What does \ur{راضی} mean? (\textquotedblleft{}Willing; agreed\textquotedblright{}.) Say it once more.
\end{tcolorbox}
\section[\texorpdfstring{\ur{اختلاف} (ikhtilāf)}{اختلاف (ikhtilāf)}]{\ur{اختلاف} (ikhtilāf) — a disagreement; a difference of opinion}

\label{lesson:UR-C530-ikhtilaf}



\begin{quote}
Before the new one: say the Urdu for one's own choice, then the Urdu for willing.
\end{quote}

\subsection*{You'll want to know: \ur{اختلاف}}
\textbf{\ur{اختلاف}} — \emph{ikhtilāf} — \textquotedblleft{}a disagreement; a difference of opinion\textquotedblright{}.

\begin{quote}
\ur{اگر} \ur{کسی} \ur{کو} \ur{اختلاف} \ur{ہے} \ur{تو} \ur{ابھی} \ur{بتائیے۔}
\end{quote}

\begin{quote}
\emph{agar kisī ko ikhtilāf hai to abhī batāiye.}
\end{quote}

\begin{quote}
If anyone disagrees, say so now.
\end{quote}

Masculine.

\begin{cousinweb}
\textbf{\ur{اختلاف}} (\emph{ikhtilāf}) is usually traced, through Persian, to Arabic \emph{ikhtilāf}, \textquotedblleft{}differing, disagreement\textquotedblright{}.
\end{cousinweb}

\begin{grammarlens}[title={What you've built}]
One more word for saying what happens if.
\end{grammarlens}

\subsection*{Your turn}
\begin{itemize}
\raggedright
  \item \emph{Say it:} \emph{ikhtilāf}
  \item \emph{Say it:} \emph{ikhtilāf}, once more
  \item \emph{Say it:} \emph{agar kisī ko ikhtilāf hai to abhī batāiye.} — If anyone disagrees, say so now.
  \item \emph{Recall:} say the Urdu for one's own choice, then the Urdu for willing, then say \emph{ikhtilāf} again
\end{itemize}

\begin{tcolorbox}[breakable,colback=teal!4,colframe=teal!35!black,title={Before you move on}]
What does \ur{اختلاف} mean? (\textquotedblleft{}A disagreement; a difference of opinion\textquotedblright{}.) Say it once more.
\end{tcolorbox}
\section[\texorpdfstring{\ur{مسترد} (mustarad)}{مسترد (mustarad)}]{\ur{مسترد} (mustarad) — rejected; turned down}

\label{lesson:UR-C530-mustarad}



\begin{quote}
Before the new one: say the Urdu for willing, then the Urdu for a disagreement.
\end{quote}

\subsection*{You'll want to know: \ur{مسترد}}
\textbf{\ur{مسترد}} — \emph{mustarad} — \textquotedblleft{}rejected; turned down\textquotedblright{}.

\begin{quote}
\ur{اگر} \ur{یہ} \ur{خط} \ur{مسترد} \ur{ہوا} \ur{تو} \ur{پھر} \ur{لکھوں} \ur{گا۔}
\end{quote}

\begin{quote}
\emph{agar yih khat mustarad huā to phir likhūṅ gā.}
\end{quote}

\begin{quote}
If this letter is turned down, I'll write again.
\end{quote}

\begin{cousinweb}
\textbf{\ur{مسترد}} (\emph{mustarad}) is usually traced, through Persian, to Arabic \emph{mustaradd}, \textquotedblleft{}taken back, returned\textquotedblright{}.
\end{cousinweb}

\begin{grammarlens}[title={What you've built}]
One more word for saying what happens if.
\end{grammarlens}

\subsection*{Your turn}
\begin{itemize}
\raggedright
  \item \emph{Say it:} \emph{mustarad}
  \item \emph{Say it:} \emph{mustarad}, once more
  \item \emph{Say it:} \emph{agar yih khat mustarad huā to phir likhūṅ gā.} — If this letter is turned down, I'll write again.
  \item \emph{Recall:} say the Urdu for willing, then the Urdu for a disagreement, then say \emph{mustarad} again
\end{itemize}

\begin{tcolorbox}[breakable,colback=teal!4,colframe=teal!35!black,title={Before you move on}]
What does \ur{مسترد} mean? (\textquotedblleft{}Rejected; turned down\textquotedblright{}.) Say it once more.
\end{tcolorbox}
\section[\texorpdfstring{\ur{پلٹنا} (palaṭnā)}{پلٹنا (palaṭnā)}]{\ur{پلٹنا} (palaṭnā) — to turn back; to go back (on one's word)}

\label{lesson:UR-C530-palatna}



\begin{quote}
Before the new one: say the Urdu for a disagreement, then the Urdu for turned down.
\end{quote}

\subsection*{You'll want to know: \ur{پلٹنا}}
\textbf{\ur{پلٹنا}} — \emph{palaṭnā} — \textquotedblleft{}to turn back; to go back (on one's word)\textquotedblright{}.

\begin{quote}
\ur{اگر} \ur{وہ} \ur{اپنی} \ur{بات} \ur{سے} \ur{پلٹا} \ur{تو} \ur{کیا} \ur{ہو} \ur{گا؟}
\end{quote}

\begin{quote}
\emph{agar voh apnī bāt se palṭā to kyā ho gā?}
\end{quote}

\begin{quote}
If he goes back on his word, what will happen?
\end{quote}

Past (he …): \ur{وہ} \ur{اپنی} \ur{بات} \ur{سے} \ur{پلٹ} \ur{گیا}, \emph{voh apnī bāt se palaṭ gayā}.

\begin{cousinweb}
\textbf{\ur{پلٹنا}} (\emph{palaṭnā}) is inherited: CDIAL lists the Hindi cousin \emph{palaṭnā}, \textquotedblleft{}to overturn\textquotedblright{}, under \emph{*pallaṭṭ}, \textquotedblleft{}turn\textquotedblright{}, a form it reconstructs (CDIAL 7968).
\end{cousinweb}

\begin{grammarlens}[title={What you've built}]
One more word for saying what happens if.
\end{grammarlens}

\subsection*{Your turn}
\begin{itemize}
\raggedright
  \item \emph{Say it:} \emph{palaṭnā}
  \item \emph{Say it:} \emph{palaṭnā}, once more
  \item \emph{Say it:} \emph{agar voh apnī bāt se palṭā to kyā ho gā?} — If he goes back on his word, what will happen?
  \item \emph{Recall:} say the Urdu for a disagreement, then the Urdu for turned down, then say \emph{palaṭnā} again
\end{itemize}

\begin{tcolorbox}[breakable,colback=teal!4,colframe=teal!35!black,title={Before you move on}]
What does \ur{پلٹنا} mean? (\textquotedblleft{}To turn back; to go back (on one's word)\textquotedblright{}.) Say it once more.
\end{tcolorbox}
